\section{Service objectives and incident learning}\label{sec:C32:06}
An operational objective needs a denominator, boundary, and observation window. Examples include the fraction of admitted tasks that complete under the specified verifier and deadline, the probability of a duplicate external effect per submitted operation, and the time from a detected compatibility failure to containment. Aggregate availability can hide a system that responds quickly but returns unusable artifacts. Report task success and effect correctness separately from endpoint availability.

Mean recovery time is sensitive to incident definition and censoring. A system that labels unresolved incidents as abandoned can appear faster. Time-to-recovery distributions should retain unresolved cases and distinguish restoration of service from reconciliation of residual effects. The latter may continue after the model endpoint is healthy. Cost accounting should include repeated model calls, tool requests, memory repair, reindexing, and human intervention according to the declared boundary.

Incident learning creates another feedback channel. A failure trace can reveal a benchmark answer, a confidential record, or an attacker payload. Automatically adding the trace to a shared memory or training corpus can contaminate future evaluation or persist the attack. Curate incident-derived records with provenance, scope, authorization, and a statement of the tested lesson. Preserve the raw restricted evidence separately when needed for audit. A procedure extracted from one incident should identify the conditions that made it appropriate.

The AgentOps research proposal provides a vocabulary for observing and optimizing agent systems, but its historical publication date and conceptual role must be distinguished from current 2026 platform capabilities~\cite{corpus250711277}. Current SDK and managed-agent documentation can establish available tracing, state, and recovery interfaces; they do not independently establish reliability in a particular workload~\cite{openai2026agents}. An operational assurance claim requires evidence from that workload and the actual release manifest.

